% !TeX root = ./../../main.tex
% Schemat blokowy procedury kalibracji i korekty metodą PDU (rysunek fig:pdu:flowchart).
% Rysunek własny, narysowany na nowo w TikZ na potrzeby rozprawy.
% Treść odpowiada algorytmom alg:pdu:calibration i alg:pdu:correction z chapters/pdu.tex.
% Wymaga bibliotek TikZ: arrows.meta, calc, positioning (wczytane w config/my-settings.tex).
\begin{tikzpicture}[
  node distance=5mm and 10mm,
  nohyph/.style={execute at begin node={\hyphenpenalty=10000\relax\exhyphenpenalty=10000\relax}},
  box/.style={draw, rounded corners=2pt, align=center, font=\footnotesize, nohyph, minimum height=10mm, text width=52mm, inner sep=3pt},
  qbox/.style={box, fill=black!10},
  cbox/.style={box, fill=white},
  arr/.style={-{Stealth[length=2mm]}, semithick},
  lab/.style={font=\footnotesize\itshape, inner sep=1.5pt},
]
% --- faza kalibracji (lewa kolumna) ---
\node[qbox] (phys) {Kalibracja fizyczna urządzenia\\(wykonuje dostawca)};
\node[qbox, below=of phys] (run) {Wykonanie $2^{m}$ obwodów referencyjnych\\$\compcirc_{1,0}(D)$, każdy $\nshots$ razy};
\node[cbox, below=of run] (err) {Wyznaczenie błędów\\$\pdu_D=\reffunexp(x_D)-\reffun(x_D)$};
\node[cbox, below=of err] (int) {Interpolacja: funkcja PDU\\$\pdu=\mathcal{I}[\mathcal{V}_m]$, zapis chwili $t_0$};
% --- faza korekty (prawa kolumna) ---
% Kolejność jak w algorytmie alg:pdu:correction: najpierw sprawdzenie warunku ponownej kalibracji, potem korekta.
\node[qbox, right=18mm of phys] (comp) {Wykonanie obwodu obliczeniowego:\\wynik $\tilde{f}$, znana faza $\varphi$};
\coordinate (dc) at ($(comp.south)+(0,-18.5mm)$);
\draw[fill=white] ($(dc)+(-29mm,0)$) coordinate (dW) -- ($(dc)+(0,14.5mm)$) coordinate (dN) -- ($(dc)+(29mm,0)$) coordinate (dE) -- ($(dc)+(0,-14.5mm)$) coordinate (dS) -- cycle;
\node[align=center, font=\footnotesize, inner sep=0pt] at (dc) {Nowa kalibracja\\fizyczna, zmiana\\układu kubitów\\lub transpilacji?};
\node[cbox, anchor=north] (corr) at ($(dS)+(0,-4mm)$) {Korekta po stronie klasycznej:\\$\hat{f}=\tilde{f}-\pdu(\varphi)$};
\coordinate (low) at ($(corr.south)+(0,-3.5mm)$);
\coordinate (rw) at ($(run.west)+(-6mm,0)$);
% --- przepływ ---
\draw[arr] (phys) -- (run);
\draw[arr] (run) -- (err);
\draw[arr] (err) -- (int);
\draw[arr] (int.east) -- ++(9mm,0) |- (comp.west);
\draw[arr] (comp.south) -- (dN);
\draw[arr] (dS) -- (corr.north);
\node[lab, anchor=north west] at ($(dS)+(0.5mm,0)$) {nie};
\draw[arr] (corr.east) -- ++(4mm,0) |- (comp.east);
\draw[arr] (dW) -- ++(-4mm,0) coordinate (t1) -- (t1 |- low) -- (rw |- low) -- (rw) -- (run.west);
\node[lab, anchor=south east] at ($(dW)+(0,0.5mm)$) {tak};
% --- nazwy faz ---
\node[lab, anchor=south] at ($(phys.north)+(0,1.5mm)$) {faza kalibracji};
\node[lab, anchor=south] at ($(comp.north)+(0,1.5mm)$) {faza korekty};
% --- legenda ---
\node[qbox, text width=6mm, minimum height=4mm, anchor=north west] (lq) at ($(phys.south west)+(0,-57mm)$) {};
\node[font=\footnotesize, anchor=west] at (lq.east) {urządzenie kwantowe};
\node[cbox, text width=6mm, minimum height=4mm, anchor=west] (lc) at ($(lq.east)+(36mm,0)$) {};
\node[font=\footnotesize, anchor=west] at (lc.east) {komputer klasyczny};
\end{tikzpicture}
